\textbf{结论名称：} 线面角的向量公式  
\textbf{适用条件：} 直线 \(l\) 与平面 \(\alpha\) 相交（若 \(l\parallel\alpha\) 则 \(\theta=0^\circ\) 仍适用）；需要取定平面的法向量 \(\vec{n}\) 与直线的方向向量 \(\vec{l}\)，且 \(\vec{n}\neq\vec{0}\)，\(\vec{l}\neq\vec{0}\)。  
\textbf{变量说明：} \(\theta\) 为直线 \(l\) 与平面 \(\alpha\) 所成角（锐角或零角），\(\vec{n}\) 为平面 \(\alpha\) 的一个法向量，\(\vec{l}\) 为直线 \(l\) 的一个方向向量。  
\textbf{核心公式：}  
\[
\boxed{\sin\theta = |\cos\langle \vec{n}, \vec{l} \rangle| = \frac{|\vec{n} \cdot \vec{l}|}{|\vec{n}| \, |\vec{l}|}}
\]  
\textbf{使用场景：} 在空间直角坐标系中求直线与平面的夹角，避免直接找射影。  
\textbf{简单例子：} 正方体 \(ABCD\)-\(A_1B_1C_1D_1\) 棱长为1，求体对角线 \(A_1C\) 与底面 \(ABCD\) 所成角。  
取 \(D(0,0,0)\)，\(C(1,0,0)\)，\(A(0,1,0)\)，\(A_1(0,1,1)\)。  
平面 \(ABCD\) 的法向量 \(\vec{n}=(0,0,1)\)，\(\vec{A_1C}=(1,-1,-1)\)。  
\[
\sin\theta=\frac{|\vec{n}\cdot\vec{A_1C}|}{|\vec{n}|\,|\vec{A_1C}|}=\frac{1}{\sqrt{3}},\quad \theta=\arcsin\frac{\sqrt{3}}{3}.
\]